\section{The seven equations encoded by the E663 packet}
All variables in each equation are universally quantified independently.
The only operation symbols are the nine in the declared signature.
The first six equations are:
\begin{align*}
N(N(x,N(N(y,x),x)),N(y,N(z,x))) &= y\qquad\text{(E1)}\\
V(P(x,s),d) &= x\qquad\text{(E2)}\\
V(P(x,s),Q(j)) &= V(s,j)\qquad\text{(E3)}\\
p_\varepsilon(s,P(a,b)) &= N(p_\varepsilon(s,a),p_\varepsilon(s,b))\qquad\text{(E4)}\\
p_\varepsilon(s,A(i,j)) &= p_E(V(s,i),V(s,j))\qquad\text{(E5)}\\
p_\varepsilon(s,Q(a)) &= N(p_\varepsilon(P(W(s,a),s),a),p_\varepsilon(P(z,s),a))\qquad\text{(E6)}
\end{align*}
For the seventh equation, use the following term abbreviations.
They are purely typographical and introduce no extra operation
symbols or axioms:
\begin{align*}
m &:= V(b,d)\\
t &:= P(b,z)\\
u &:= Q(b)\\
e &:= p_\varepsilon(P(t,P(u,b)),m)\\
f &:= N(p_E(A(b,z),t),N(d,p_E(A(b,z),z)))\\
g &:= N(d,p_E(y,u))\\
h &:= p_\varepsilon(P(y,P(x,b)),m)\\
j &:= N(p_E(W(x,d),x),N(d,p_E(x,u)))\\
k &:= N(p_E(y,u),p_\varepsilon(P(y,P(P(b,y),b)),m))\\
l &:= N(N(g,N(h,j)),k)
\end{align*}
The last equation is
\[d=N(N(e,N(d,f)),N(d,l)).\tag{E7}\]
The ground query is $Q(d)=d$. The appendix generator expands every
abbreviation and checks exact equality with the seventh JSON term
tree, so the displayed presentation matches the encoded packet.
The inherited calibration and the reasons for the noncollapse query
are retained in the pinned 235-state archive.
